% !TeX root = ../泛函分析笔记.tex
% 由主文档通过 \input 载入；请编译主文档。
\setcounter{chapter}{16} % 保留原书章号。
\chapter{算子半群与演化方程}
\label{chap:fa-semigroups}

本章为选读内容。有限维方程 $u'(t)=Au(t)$ 的解由矩阵指数给出。无限维中，即使生成时间演化的算子 $A$ 无界，每个固定时刻的演化仍可能由有界算子描述。

\section{强连续半群}
\begin{definition}
Banach 空间 $X$ 上的\emph{强连续半群（strongly continuous semigroup）}是算子族 $S(t)\in\mathcal B(X)$，$t\geq0$，满足
\[
 S(0)=I,\qquad S(t+s)=S(t)S(s),
\]
并且对每个 $x\in X$，映射 $t\mapsto S(t)x$ 连续。若 $\|S(t)\|\leq1$，称为压缩半群。
\end{definition}
当 $A\in\mathcal B(X)$ 时，$S(t)=\sum_{n\geq0}t^nA^n/n!$ 按算子范数收敛；绝对收敛允许重排双重级数，从二项式公式得到半群性质，而且 $S'(t)=AS(t)$。无界情形不应直接写未经定义的算子幂级数。

任意强连续半群在有限时间区间上算子范数一致有界：对每个 $x$，轨道在该区间上有界，再用一致有界原理。取 $M=\max\{1,\sup_{0\leq t\leq1}\|S(t)\|\}$，由半群性质得到 $\|S(t)\|\leq M e^{\omega t}$，其中 $\omega=\log M$。

“半群”只要求非负时间，是因为演化未必能够稳定地倒退。矩阵指数 $e^{tA}$ 在有限维中总有逆 $e^{-tA}$；热方程的正向演化不断压低高频误差，倒退则要把这些坐标指数放大，通常不能得到整个 $L^2$ 空间上的有界逆。因此半群性质描述时间拼接，并不自动提供负时间。

\section{生成元与闭性}
\begin{definition}
半群的\emph{无穷小生成元（infinitesimal generator）}由
\begin{equation}
 Ax=\lim_{h\downarrow0}\frac{S(h)x-x}{h},\qquad
 D(A)=\left\{x\in X:\text{此极限在 }X\text{ 中存在}\right\}
 \label{eq:fa-semigroup-generator}
\end{equation}
定义。
\end{definition}
若 $x\in D(A)$，由半群性质可得 $S(t)x\in D(A)$，且 $AS(t)x=S(t)Ax$；右、左差商进一步给出 $\frac{\mathrm d}{\mathrm dt}S(t)x=S(t)Ax$（在 $t=0$ 取右导数）。因此
\begin{equation}
 S(t)x-x=\int_0^tS(s)Ax\,\mathrm ds.
 \label{eq:fa-generator-integral}
\end{equation}
连续 Banach 空间值函数的积分可定义为 Riemann 和的范数极限；一致连续性和完备性保证极限存在。

生成元的定义只使用固定初值 $x$ 的差商。在有界情形可以把极限统一为算子范数极限，得到处处定义的 $A$；在无界情形，某些初值的轨道只有连续性，没有初始导数。定义域 $D(A)$ 因而可以理解为允许在初始时刻写出微分方程的那一部分初值，而不是人为附加的技术集合。

\begin{proposition}
强连续半群的生成元稠密定义且闭。
\end{proposition}
\begin{proof}
对任意 $x\in X$，令 $x_h=h^{-1}\int_0^hS(s)x\,\mathrm ds$。利用半群性质移动积分区间可得 $x_h\in D(A)$ 且 $Ax_h=(S(h)x-x)/h$；强连续性给出 $x_h\to x$，故定义域稠密。

若 $x_n\in D(A)$、$x_n\to x$、$Ax_n\to y$，在式\eqref{eq:fa-generator-integral}中利用有限时间区间上的统一算子界取极限，得到 $S(t)x-x=\int_0^tS(s)y\,\mathrm ds$。除以 $t$ 并令 $t\downarrow0$，右侧趋于 $y$，所以 $x\in D(A)$ 且 $Ax=y$，即 $A$ 闭。
\end{proof}

\section{预解算子与 Hille--Yosida 定理}
对闭算子 $A$，$\lambda\in\rho(A)$ 指 $\lambda I-A:D(A)\to X$ 双射且其逆作为 $X\to X$ 的算子有界。若 $A$ 生成满足 $\|S(t)\|\leq Me^{\omega t}$ 的半群，则对实数 $\lambda>\omega$，
\[
 R(\lambda,A)x=\int_0^\infty e^{-\lambda t}S(t)x\,\mathrm dt,
 \qquad\|R(\lambda,A)\|\leq\frac M{\lambda-\omega}.
\]
积分由增长估计绝对收敛。其逆关系可先在 $D(A)$ 中用式\eqref{eq:fa-generator-integral}分部积分验证，再利用稠密性和生成元闭性扩展；这是半群与预解估计之间的联系。

\begin{theorem}[Hille--Yosida 的压缩半群版本；引用]
设 $A$ 是 Banach 空间上的稠密定义闭算子。它生成强连续压缩半群，当且仅当每个实数 $\lambda>0$ 属于 $\rho(A)$，且 $\|R(\lambda,A)\|\leq1/\lambda$。
\end{theorem}
必要性由上述积分公式得到。充分性通过有界近似 $A_\lambda=\lambda A R(\lambda,A)=\lambda^2R(\lambda,A)-\lambda I$ 及指数半群的极限构造证明；本书不展开这一极限论证，完整证明见\cite{engel-nagel}第二章第三节。一般增长常数 $M$ 的版本需要对预解算子的各次幂给出相应估计，不能只替换单次估计中的常数。

\section{经典解与温和解}
设 $A$ 生成 $S(t)$。若 $u_0\in D(A)$，则 $u(t)=S(t)u_0$ 是 $u'=Au$ 的经典解，即轨道可微、取值在 $D(A)$ 中并满足微分方程。唯一性可固定终点 $t$，对 $s\mapsto S(t-s)u(s)$ 求导，得到导数为零。

一般 $u_0\in X$ 未必属于定义域，仍可定义 $u(t)=S(t)u_0$，称为\emph{温和解（mild solution）}。对连续外力 $g:[0,T]\to X$，非齐次问题的温和解为
\[
 u(t)=S(t)u_0+\int_0^tS(t-s)g(s)\,\mathrm ds.
\]
连续外力保证这条公式有意义，但要逐项微分并取得经典解，还需初值与外力的额外正则性；不能把温和解公式本身当作经典可微性的证明。

\section{区间上的热半群}
采用例\ref{ex:fa-green}中的正交基 $e_n(x)=\sqrt2\sin(n\pi x)$。若 $u_0=\sum_nc_ne_n\in L^2(0,1)$，令
\[
 S(t)u_0=\sum_{n=1}^{\infty}e^{-(n\pi)^2t}c_ne_n.
\]
系数逐项满足乘法规律，故这是半群；Parseval 恒等式与级数的支配收敛给出强连续性，且 $\|S(t)\|=e^{-\pi^2t}$。其生成元为
\[
 Au=-\sum_n(n\pi)^2c_ne_n,\qquad
 D(A)=\left\{u:\sum_n(n\pi)^4|c_n|^2<\infty\right\}.
\]
定义域内，差商由 $|(e^{-ah}-1)/h|\leq a$ 控制，故按 $L^2$ 范数收敛；反之，若差商收敛，逐坐标极限迫使所得系数平方可和，恰好给出该定义域。

定义 $H^2(0,1)=\{u\in H^1(0,1):u''\in L^2(0,1)\}$，二阶导数仍为弱导数。由 Green 算子的逆关系可将上述定义域识别为 $H^2(0,1)\cap H_0^1(0,1)$，并有 $Au=u''$。所以 $u(t)=S(t)u_0$ 描述零边界条件下的热方程 $u_t=u_{xx}$。每个 $t>0$ 的指数衰减还保证 $u(t)\in D(A)$；在初始时刻是否可微则取决于 $u_0$ 是否属于 $D(A)$。

\begin{example}[两个频率的衰减速度]
若 $u_0=e_1+2e_3$，则
\[
 u(t)=e^{-\pi^2t}e_1+2e^{-9\pi^2t}e_3,
 \qquad\|u(t)\|_2^2=e^{-2\pi^2t}+4e^{-18\pi^2t}.
\]
\end{example}
\begin{solve}
第一式逐项应用半群的定义，第二式由正交性得到。第三频率的衰减指数是第一频率的九倍，所以较晚时刻主要剩下第一项。在无穷多个频率同时存在时，对每个固定初值，尾部系数的平方和已经很小，故 $S(t)u_0\to u_0$；但对每个 $t>0$，总能找到足够大的 $n$，使 $1-e^{-(n\pi)^2t}$ 任意接近一，所以 $\|S(t)-I\|=1$。强连续与算子范数连续在这个例子中的差别，来自是否允许测试向量随时间改变。
\end{solve}

更定量地，对固定 $t>0$，$\sup_{r\geq0}r^2e^{-2rt}<\infty$，取 $r=(n\pi)^2$ 得
\[
 \sum_n(n\pi)^4e^{-2(n\pi)^2t}|c_n|^2
 \leq\left(\sup_{r\geq0}r^2e^{-2rt}\right)\sum_n|c_n|^2<\infty.
\]
这证明任意 $L^2$ 初值在正时间后进入生成元定义域。相反，若要从时刻 $t$ 的系数倒推初值，需要乘以 $e^{(n\pi)^2t}$，这些放大倍数没有统一上界，不能得到有界的反向演化算子。

\section{习题}
\begin{enumerate}
 \item 对有界 $A$ 证明 $\|e^{tA}\|\leq e^{t\|A\|}$。
 \item 在上述热半群中证明 $\|S(t)-I\|=1$ 对每个 $t>0$ 成立，说明强连续不等于算子范数连续。
 \item 设 $u_0\perp e_1,\ldots,e_{m-1}$，证明 $\|S(t)u_0\|_2\leq e^{-(m\pi)^2t}\|u_0\|_2$，并说明常数的最优性。
 \item 证明非负自伴算子的谱演算产生强连续压缩半群 $e^{-tA}$。
\end{enumerate}

